\documentclass[10pt,a4paper]{amsart}
\usepackage[margin=19mm]{geometry}
\usepackage{amsmath,amssymb,mathtools}
\usepackage[scr=rsfs,cal=euler]{mathalfa}
\usepackage{xfrac}
\usepackage[unicode,hidelinks,bookmarks=false]{hyperref}
\newcommand{\ie}{i.e.\ }
\newcommand{\cf}{cf.\ }
\newcommand{\resp}{respectively\ }
\newcommand{\Subsection}{\subsection}
\providecommand{\nobreakdash}{\nobreak\hbox{-}}
\setlength{\emergencystretch}{2em}
\multlinegap0pt
\usepackage{amssymb}
\usepackage[shortlabels]{enumitem}
\usepackage{xcolor}
\usepackage{tikz}
\usepackage{float}
\usepackage{mathtools}
\newtheorem{lemma}{Lemma}
\usepackage{cleveref}
\crefname{lemma}{Lemma}{Lemmas}
\title{On the number of exceptional intervals to the prime number theorem in short intervals}
\author{Ayla Gafni, Terence Tao}
\date{Source: arXiv:2505.24017v1 (CC BY 4.0)}
\begin{document}
\maketitle

\noindent\emph{Source and licence.} On the number of exceptional intervals to the prime number theorem in short intervals. Version arXiv:2505.24017v1 (CC BY 4.0), distributed by arXiv under the Creative Commons Attribution 4.0 International licence, http://creativecommons.org/licenses/by/4.0/, as recorded in the arXiv licence metadata for this exact version and shown on the abstract page https://arxiv.org/abs/2505.24017v1. Full licence notice: \texttt{licenses/arxiv-2505.24017-CC-BY-4.0.txt}.\par\smallskip

\noindent\textit{Editorial scope.} Counted unit: the article's lemma linfty-bound (source0.tex lines 418--421). The article's complete original proof of it is delivered with it (source0.tex lines 423--429, one \texttt{proof} environment of 7 lines). No mathematics is written, repaired or re-derived: the statement and the proof are the article's own lines, in the article's own notation, and no retained line is substituted or re-flowed.  Retained uncounted as the source-local context the proof needs: lines 98--100; lines 390--395.  Nothing was omitted from any retained span, so the package carries no omission record.  No retained line contains a citation, so the wrapper carries no bibliography and no bibliography entry.  No retained line contains a figure environment, an image-inclusion command or drawing code, so no image is delivered and the package declares no asset dependency.  Editorial formatting only, in addition: an AMS \texttt{amsart} preamble that restates the article's own declarations and macros used by the retained lines, namely the environments \texttt{lemma} on separate counters, together with the standard packages those lines need.  The retained environments print in this document as Lemma 1 in order of appearance, whereas the article prints the same items with the numbering of its own full document.  The complete native source of the article as distributed in the arXiv e-print for this exact version is bundled unchanged as \texttt{sources/179022/source0.tex}.
\begin{equation}\label{nsig-up}
  N(\sigma, T) \leq T^{A(\sigma)(1-\sigma) + o(1)}
\end{equation}

\begin{lemma}[Contribution near the right edge]\label{right-edge}  If $\eta_0>0$ is sufficiently small depending on $\theta$, and $I \subset [1-\eta_0,1]$, then
  \begin{equation}\label{si-bound}
\sup_{X \leq x \leq 2X}|S_I(x)| \ll_\theta \exp\left(-c_\theta \frac{\log^{1/3} X}{\log\log^{1/3} X} \right)\frac{X}{\tau}
  \end{equation}
  for some $c_\theta>0$ depending only on $\theta$.
\end{lemma}

\begin{lemma}[$L^\infty$ bound]\label{linfty-bound}  If $I \subset [\sigma_-, \sigma_+]$ for some fixed $0 \leq \sigma_- \leq \sigma_+ \leq 1$, then
$$ \sup_{X \leq x \leq 2X} |S_I(x)| \ll_\theta X^{(1-\theta)(1-\sigma_-) A(\sigma_-) + \theta + \sigma_+ - 1 + o(1)}$$
as $X \to \infty$.
\end{lemma}

\begin{proof}
By repeating the arguments used to prove \Cref{right-edge}, we have
$$ S_I(x) \ll \frac{X \log X}{\tau} \sup_{\sigma_- \leq \sigma \leq \sigma_+} N(\sigma,T) X^{\sigma-1}.$$
If we crudely bound $N(\sigma,T) \leq N(\sigma_-,T)$ and $X^{\sigma-1} \leq X^{\sigma_+ - 1}$, we obtain
$$ S_I(x) \ll \frac{X \log X}{\tau} N(\sigma_-,T) X^{\sigma_+ - 1}.$$
The claim now follows from \eqref{nsig-up} and the definitions of $T,\tau$.
\end{proof}

\end{document}
